\section{INTRODUCCIÓN}

\subsection{Antecedentes}
\dots

El virus sincitial respiratorio (RSV) plantea otro problema de vigilancia por su carga en lactantes, adultos mayores e inmunocomprometidos, y por su presentación clínica compartida con otros virus respiratorios. Hughes et al.(2022) detectaron ARN de RSV en sólidos sedimentados de aguas residuales de dos plantas y reportaron asociación fuerte con tasas de positividad clínica (Kendall tau = 0,65--0,77; $p < 10^{-7}$) \cite{Hughes_2022}. Le et al. (2025) desarrollaron un flujo de secuenciación ONT para muestras clínicas y aguas residuales, obtuvieron genomas completos en el 98 \% de las muestras con más de 10.000 copias/mL y observaron que las frecuencias de clados en aguas residuales representaban las frecuencias observadas en pacientes \cite{Le_2025}. El mismo estudio subraya la necesidad de una vigilancia genómica accesible en contextos con recursos limitados, debido a la menor representación de datos genómicos de RSV procedentes de países de ingresos bajos y medios \cite{Le_2025}.

En Quito, la PTAR Quitumbe es un espacio estratégico para observar el comportamiento de esta brecha. En el trabajo de Mejía Calle (2024) se realizó un muestreo pasivo semanal de 24 horas en el influente de la PTAR Quitumbe. En esa investigación, los genes N1 y N2 de SARS-CoV-2 fueron detectados en todas las muestras; en la secuenciación se utilizó tecnología de Oxford Nanopore permitiendo identificar linajes de Ómicron, incluida la detección de JN.1 en aguas residuales antes de su identificación clínica local \cite{MejiaCalle_2024}. Esta evidencia confirma la factibilidad del sitio para la vigilancia ambiental, que se centra en SARS-CoV-2.

\citeA{Sims_2020} situaron la epidemiología basada en aguas residuales como una herramienta para monitorear la propagación de enfermedades infecciosas a escala comunitaria; luego, en el estudio de \cite{Parkins_2023} la consolidaron como estrategia para vigilancia tras la pandemia de COVID-19. Sobre esa base, las secciones siguientes abordan la epidemiología basada en aguas residuales, la biología de SARS-CoV-2, RSV y poliovirus, las técnicas de detección molecular, la secuenciación y el análisis bioinformático de las señales virales aplicadadas en una matriz ambiental compleja.

\subsection{Vigilancia basada en aguas residuales}

La vigilancia basada en aguas residuales analiza de forma sistemática muestras de alcantarillado o de influentes de plantas de tratamiento para identificar marcadores biológicos asociados con una población. El punto de muestreo define qué grupo de la población está siendo representado, además del alcance de la vigilancia y el tipo de la información que puede obtenerse para la toma de decisiones \cite{Parkins_2023}.

... 

En aguas residuales, la caracterización genómica revela qué variantes presentan una señal positiva de naturaleza poblacional, ya que una misma muestra contiene genomas de múltiples individuos y de distintos estados de infección. Su valor epidemiológico depende de criterios explícitos de cobertura, profundidad, calidad de lectura y abundancia relativa \cite{Parkins_2023,Child_2023}. La representación de los linajes minoritarios puede alterarse por procesos técnicos como la degradación del ARN, la competencia entre amplicones y la elección de la referencia, factores que \citeA{MejiaCalle_2024} y \citeA{Le_2025} reconocieron al trabajar con muestras de baja carga.

\subsection{Herramientas de análisis bioinformático}

\subsubsection{Procesamiento de secuencias}

El procesamiento bioinformático sigue etapas ordenadas: demultiplexado, eliminación de adaptadores, filtrado de calidad, mapeo con referencias genómicas y generación de una secuencia consenso o ensamblaje, con estimación de la cobertura por posición. En la secuenciación de ONT, el flujo añade el basecalling y el control de calidad de lecturas largas, por otro lado en plataformas de secuenciación como Illumina se prioriza la calidad por base y el manejo de lecturas pareadas \cite{Dash_2024}.

Las muestras de aguas residuales contienen mezcla de material biológico, sólidos, inhibidores residuales y ácidos nucleicos degradados,  situación que incrementa su complejidad de análisis genómico. \citeA{Child_2023} mostraron que la elección entre metagenómica, captura híbrida y PCR dirigida modifica la cobertura recuperada de virus humanos. En la PTAR Quitumbe, se estimó los linajes de SARS-CoV-2 usando Freyja a partir de las lecturas ambientales, con un rendimiento que dependió de la carga viral y de la calidad de amplificación realizada en la PCR \cite{MejiaCalle_2024}. Es importante destacar que la trazabilidad del análisis sostiene su reproducibilidad, cada muestra debe conservar metadatos de fecha, punto de muestreo, método de concentración, controles, valores de Cq, plataforma y versión de las herramientas, sin los cuales la comparación temporal o entre patógenos resulta menos confiable \cite{Parkins_2023}.

\subsubsection{Caracterización genética}

La caracterización genética de una entidad infecciosa utiliza herramientas de asignación de linajes, comparando lecturas o consensos obtenidos en la secuenciación, con bases de datos de referencia actualizadas. Freyja es una herramienta bioinformática que estima proporciones de variantes de un virus en muestras mixtas \cite{Karthikeyan_2022}. En la PTAR Quitumbe, se usó este software para asignar linajes de Ómicron a partir de datos ambientales \cite{MejiaCalle_2024}. \citeA{Crits_Christoph_2021} identificaron variantes prevalentes en aguas residuales de la región, cuando la cobertura era suficiente en las posiciones informativas del genoma del agente infeccioso; es decir, sitios que presentan mutaciones características que permiten distinguir una variante de otra.

En RSV, la caracterización genética distingue RSV-A y RSV-B y selecciona referencias compatibles con el subgrupo predominante; por ejemplo, \citeA{Le_2025} desarrollaron RSVTyper, que es una herramienta que genera la detección automática del subgrupo y la selección de referencia del virus. Por otro lado, la identificación de poliovirus se apoya en la diferenciación intratípica y en la secuenciación de la región VP1 para reconocer cepas relacionadas o derivadas de vacuna \cite{Asghar_2014,Klapsa_2022}. En muestras ambientales donde hay mezcla de genomas, la detección de una variante minoritaria se logra mediante una cobertura del genoma o de las regiones diagnósticas, concordancia entre réplicas y frecuencia relativa del linaje; la finalidad es reducir los reportes engañosos por error de secuenciación o contaminación cruzada \cite{Dash_2024,Child_2023}.

\subsubsection{Interpretación de resultados y limitaciones}

Una señal positiva indica presencia de material genético compatible con el blanco analizado, pero su magnitud varía por los determinantes de la matriz utilizada, de modo que las tendencias y las comparaciones de resultados de las muestras obtenidas bajo condiciones similares son más informativas que las lecturas aisladas \cite{Parkins_2023,Hughes_2022}.

Las limitaciones genómicas se centran en la baja concentración viral, el ARN fragmentado, el predominio de material bacteriano y los sesgos de amplificación. \citeA{Child_2023} documentaron que la metagenómica shotgun puede ser insuficiente para recuperar genomas virales sin enriquecimiento previo, y que los métodos dirigidos aumentan la cobertura del microorganismo, pero se debe tener conocimiento previo del genoma de este. En la PTAR Quitumbe, las muestras con Cq tardíos tuvieron menor éxito de secuenciación y los cambios de cebadores afectaron la asignación de linajes \cite{MejiaCalle_2024}. Una planta de tratamiento representa a una población conectada al alcantarillado, con variaciones por movilidad, cobertura de red y aportes no domésticos; los datos ambientales evidencian circulación y tendencias comunitarias, mientras la atribución a barrios o grupos específicos requiere diseños y metadatos adicionales \cite{Parkins_2023,Sims_2020}.

\subsection{Contexto epidemiológico del estudio}

\subsubsection{Situación en Ecuador}

La evidencia local de vigilancia ambiental proviene del antecedente de SARS-CoV-2 en la PTAR Quitumbe, donde \citeA{MejiaCalle_2024} identificó linajes de Ómicron mediante muestreo pasivo y posterior secuenciación usando Oxford Nanopore, incluida la detección ambiental de JN.1 antes de su identificación clínica local.

\citeA{MejiaCalle_2024} destacó que existe una carencia de información genómica clínica en Ecuador, asociada con la disponibilidad limitada de datos actualizados y la carga irregular de secuencias en los repositorios públicos. La vigilancia basada en aguas residuales ofrece en este escenario estimaciones poblacionales menos sesgadas que la prueba clínica individual y a un costo relativo menor \cite{Karthikeyan_2022}.

Para RSV y poliovirus, los recursos revisados son evidencia internacional de la vigilancia ambiental y la detección de circulación silenciosa, pero su información específica es deficiente para el contexto ecuatoriano, manifestando la necesidad de generar evidencia nacional en este campo. La PTAR Quitumbe permite ampliar la vigilancia desde SARS-CoV-2 hacia RSV y poliovirus, dos objetivos de vigilancia sanitaria documentada \cite{Asghar_2014,Hughes_2022,Bottcher_2025}.

